\documentclass[11pt,a4paper]{article}
\usepackage[margin=20mm]{geometry}
\usepackage{fontspec}
\setmainfont{texgyrepagella}[Extension=.otf,UprightFont=*-regular,BoldFont=*-bold,ItalicFont=*-italic,BoldItalicFont=*-bolditalic]
\setsansfont{texgyreheros}[Extension=.otf,UprightFont=*-regular,BoldFont=*-bold,ItalicFont=*-italic,BoldItalicFont=*-bolditalic]
\usepackage{microtype}
\usepackage[normalem]{ulem}
\usepackage{siunitx}
\usepackage{enumitem}
\DeclareSIUnit{\bp}{bp}
\DeclareSIUnit{\flops}{FLOPs}
\newcommand{\deleted}[1]{{\color{removed}\sout{#1}}}
\newcommand{\inserted}[1]{{\color{added}#1}}
\usepackage{xcolor}
\usepackage[breakable,skins]{tcolorbox}
\usepackage{fancyhdr}
\usepackage[hidelinks]{hyperref}
\definecolor{heading}{HTML}{000000}
\definecolor{muted}{HTML}{526575}
\definecolor{removed}{HTML}{982832}
\definecolor{added}{HTML}{176037}
\newtcolorbox{reviewcomment}{enhanced,breakable,sharp corners,boxrule=0pt,
  colback=gray!4,coltext=black,borderline west={1.5pt}{0pt}{muted!65},
  left=9pt,right=9pt,top=6pt,bottom=6pt,before skip=6pt,after skip=9pt}
\newtcolorbox{removedtext}{enhanced,breakable,sharp corners,boxrule=0pt,
  colback=red!4,coltext=removed,borderline west={2pt}{0pt}{removed},
  title={$-$},colbacktitle=red!4,coltitle=removed,fonttitle=\bfseries,
  left=9pt,right=9pt,top=6pt,bottom=6pt,before skip=6pt,after skip=0pt}
\newtcolorbox{addedtext}{enhanced,breakable,sharp corners,boxrule=0pt,
  colback=green!4,coltext=added,borderline west={2pt}{0pt}{added},
  title={$+$},colbacktitle=green!4,coltitle=added,fonttitle=\bfseries,
  left=9pt,right=9pt,top=6pt,bottom=6pt,before skip=0pt,after skip=10pt}
\setlength{\parindent}{0pt}
\setlength{\parskip}{7pt}
\setlength{\headheight}{14pt}
\pagestyle{fancy}
\fancyhf{}
\renewcommand{\headrulewidth}{0pt}
\fancyfoot[R]{\footnotesize\sffamily\color{muted}\thepage}
\hypersetup{pdftitle={Response to the examination jury},pdfauthor={Florian De Rop}}
\begin{document}

{\LARGE\sffamily\bfseries\color{heading}Response to the examination jury\par}

\section*{\sffamily\color{heading}Professor Céline Vallot}

\subsection*{\sffamily\color{heading}1. Critical perspective and technology selection}

\begin{reviewcomment}
Given the breadth of the field, however, some sections of the Introduction remain relatively descriptive and could have benefited from a more integrated and critical perspective. In particular, I would have appreciated further discussion of how these different technologies complement one another, and what their respective strengths and limitations are. Providing more guidance on when and why one approach may be preferred over another, depending on the biological question asked, would also help the reader navigate the rapidly evolving single-cell genomics landscape. These points could provide interesting topics for further discussion during the oral defense.
\end{reviewcomment}

I have further explored the connections between these technologies and how they complement one another. Given the length of the introduction, I have chosen to prepare notes on these points for the question-and-answer session during the oral defense rather than rewrite the introduction.

\section*{\sffamily\color{heading}Professor Jeroen Lammertyn}

\subsection*{\sffamily\color{heading}1. Integration of the general abstract}

\begin{reviewcomment}
The general abstract is simply composed of two paragraphs discussing the two papers of the PhD. I would add a few sentences on linking both topics (e.g., a few general introductory lines in the abstract) and a kind of integrated conclusion, aspects that are well worked out for instance in the Layman's Abstract on the next page.
\end{reviewcomment}

I agree that the abstract needed more introductory context and a clearer connection between the two studies. I have modified the abstract as follows:

\begin{reviewcomment}
\setlength{\parskip}{0.5\baselineskip}
\inserted{Gene expression and chromatin accessibility profiling provide complementary views of gene regulation across tissues, developmental stages, and disease. While bulk sequencing assays average these signals across populations of cells, single-cell RNA-seq and single-cell ATAC-seq can resolve cellular heterogeneity and enable the characterization of individual cell types and states. The rapid development of single-cell technologies has highlighted two related challenges: developing methods that make high-quality single-cell profiling more accessible, and understanding how differences between technologies influence the biological conclusions drawn from the resulting data. This thesis addresses both challenges through the development of an open-source single-cell sequencing platform and a systematic benchmark of scATAC-seq technologies.}

\deleted{Single-cell RNA-seq and} \inserted{To make} single-cell \deleted{ATAC-seq have enabled the building of cell type atlases across organisms, tissues, and disease. This thesis centers around the development of} \inserted{profiling more accessible, we developed} HyDrop, a flexible, open-source droplet platform for scATAC-seq and scRNA-seq based on dissolvable polyacrylamide hydrogel beads. These beads carry split-pool barcodes which index RNA or accessible chromatin from single cells. We optimized the protocol by tuning primer load, adding a clathrin inhibitor, and titrating different enzymes to increase yield. HyDrop-ATAC generated \num{7996} high-quality single-cell profiles from flash-frozen mouse cortex. HyDrop-RNA produced \num{9508} transcriptomes from the same tissue. We also recovered a small population of rare, FACS-sorted \textit{Drosophila} neurons, demonstrating applicability to low-input samples. Across mouse cortex, HyDrop-RNA reached a median \num{3389} UMIs and \num{1658} genes per cell, and HyDrop-ATAC recovered a median \num{4148} unique fragments per cell with high TSS enrichment. Chemistry, microfluidic chips, and protocols are open-source.

\deleted{We also} \inserted{To assess how technology choices affect data quality and downstream analyses, we systematically compared eight scATAC-seq technologies, including HyDrop. To enable consistent processing across technologies, we} developed \texttt{PUMATAC}, a \deleted{containerized, universal scATAC-seq} \inserted{containerized} preprocessing framework that \deleted{unifies} \inserted{standardizes} trimming, whitelist-aware barcode correction, and alignment across different scATAC-seq formats. \deleted{We used \texttt{PUMATAC} to} \inserted{The} benchmark \inserted{comprised} eight scATAC-seq methods \deleted{on} \inserted{applied to} human PBMCs: 10x v1, v1.1, v2, 10x Multiome, Bio-Rad ddSEQ, mtscATAC-seq, s3-ATAC, and HyDrop, for a total of \num{47} libraries and \num{169227} high-quality cells. We detected various systematic differences across these technologies. Sequencing efficiency, defined as fraction of reads retained in the final cell-by-peak matrix, ranged from \textasciitilde\num{4}\% for s3-ATAC and HyDrop to \textasciitilde\num{28}\% for 10x v2. Per-cell sensitivity scaled with sequencing efficiency and impacted cell-type annotation accuracy, genotype demultiplexing, peak calling, differential accessibility, and transcription factor motif enrichment. Methods also differed in specificity: fragment promoter and TSS proximity as well as aggregate fraction of fragments in peaks varied markedly, with ddSEQ recovering more TSS-proximal fragments than 10x. This effect led to technique-dependent differentially accessible regions, and imperfect cross-technology integration even after batch correction.

\inserted{Together, these studies address both the development and evaluation of single-cell sequencing technologies. HyDrop provides an accessible platform for generating single-cell RNA and chromatin accessibility profiles, while the benchmark demonstrates that technical differences between scATAC-seq methods can substantially affect sensitivity, specificity, and downstream biological analyses. These results highlight that improvements in single-cell genomics should be evaluated not only in terms of throughput, but also by the quality and biological informativeness of the signal recovered from each cell. HyDrop provides a platform on which such improvements can be developed, while the benchmarking framework provides a systematic way to assess their impact.}
\end{reviewcomment}

\subsection*{\sffamily\color{heading}2. Reference formatting}

\begin{reviewcomment}
Do you have a particular motivation to implement two strategies for referencing? A general bibliography at the end and working with footnotes in the introduction and discussion chapter?
\end{reviewcomment}

The footnotes allow readers to identify the cited papers while reading, without having to scroll to the bibliography. The general bibliography provides a complete list of references at the end of the thesis.

\subsection*{\sffamily\color{heading}3. Recent literature}

\begin{reviewcomment}
The most recent publications date back to 2025. Did the candidate screen for more recent information as well? I think this is a relevant remark in this very rapidly evolving field.
\end{reviewcomment}

In the new and extended discussion, I have included several works and projects from 2026.

\subsection*{\sffamily\color{heading}4. Personal bibliography}

\begin{reviewcomment}
Can the candidate add a personal bibliography including the conferences, journal papers, lectures, etc.?
\end{reviewcomment}

I have expanded the Author profile with my journal publications, published laboratory protocols, and conference talks and seminar presentations.

\subsection*{\sffamily\color{heading}5. Future directions, bottlenecks, and challenges}

\begin{reviewcomment}
In the discussion section it might be interesting to elaborate more on where the field should go and what the major bottlenecks and challenges are.
\end{reviewcomment}

In the new and extended discussion, I discuss what I think are major bottlenecks, and how industry and academia are dealing with these issues in their own ways.

A diff of chapter 4 is included in Appendix~\ref{app:chapter4-diff}.

\section*{\sffamily\color{heading}Professor Klaas Mulder}

Professor Mulder approved the manuscript without revisions.

\section*{\sffamily\color{heading}Professor Thierry Voet}

\subsection*{\sffamily\color{heading}1. Balance between the introduction and discussion}

\begin{reviewcomment}
There is a bit of an imbalance between the first chapter (introduction) and the last chapter (discussion). While the introduction is (at times overly) detailed, the discussion is relatively brief and succinct. Consider balancing this a bit more.
\end{reviewcomment}

I have amended the discussion with a table of ongoing large omics projects. I agree that the discussion is proportionally too short, however, due to time constraints I cannot put all my ideas to rest there. I plan to remain active in the field and release written content every now and then.

\subsection*{\sffamily\color{heading}2. Microarray probe count (p. 11)}

\begin{reviewcomment}
p.11 states that ‘While microarrays allow for up to 50 000 probes to be interrogated in a single experiment', however high-density arrays allow for many more probes. Please correct.
\end{reviewcomment}

The sentence was revised as follows.

\begin{removedtext}
While microarrays allow for up to 50 000 probes to be interrogated in a single experiment, their limited dynamic range, high false-positive hybridisation and heterogeneous probe hybridisation capacities skew absolute quantification.
\end{removedtext}
\begin{addedtext}
While high-density microarrays allow millions of probes to be interrogated in a single experiment, their limited dynamic range, high false-positive hybridisation and heterogeneous probe hybridisation capacities skew absolute quantification.
\end{addedtext}

\subsection*{\sffamily\color{heading}3. Sequencing errors and aligner choice (p. 22)}

\begin{reviewcomment}
p. 22 states "sequencing by synthesis (Illumina) generates more deletion artifacts due to incorporation errors, and these need to be explicitly corrected for by bwa-mem, whereas minimap2 is better suited for handling the higher single-base error rates found in Nanopore/PacBio reads". However, Illumina SBS generates primarily substitution artefacts, while nanopore has more deletion artefacts?
\end{reviewcomment}

The two sentences about read length and sequencing errors were replaced with the following three sentences.

\begin{removedtext}
Short-read aligners such as bowtie2 are optimized for reads < 200 bp, while long-read aligners like minimap2 are better at handling reads longer than 1000 bp.

The type of sequencing technology used also plays a role: sequencing by synthesis (Illumina) generates more deletion artifacts due to incorporation errors, and these need to be explicitly corrected for by bwa-mem, whereas minimap2 is better suited for handling the higher single-base error rates found in Nanopore/PacBio reads.
\end{removedtext}
\begin{addedtext}
For genomic alignment, bowtie2 and bwa-mem are commonly used for Illumina short reads.

Minimap2 is generally better suited to long reads from Nanopore and PacBio.

Sequencing errors differ in frequency and type between platforms, and alignment parameters need to account for these differences.
\end{addedtext}

\subsection*{\sffamily\color{heading}4. Missing space (p. 39)}

\begin{reviewcomment}
p. 39 typo : “10\%mouse cells”
\end{reviewcomment}

Corrected “10\%mouse cells” to “10\% mouse cells”.

\subsection*{\sffamily\color{heading}5. Missing space (p. 43)}

\begin{reviewcomment}
p. 43 : typo “percluster”
\end{reviewcomment}

Corrected “percluster” to “per cluster”.
\clearpage
\appendix
\section[Chapter 4 changes]{\sffamily\color{heading}Chapter 4 changes}
\label{app:chapter4-diff}
The comparison below shows changes from the September 10 version of chapter 4.

\begingroup
\sloppy
\setlength{\parskip}{0.5\baselineskip}
% Compared against 1039e4b (2026-09-10). Citation markers omitted; table replaced by placeholder.
The work described in this thesis encompasses two parts. First, we developed HyDrop, an open-source droplet microfluidic platform for scATAC-seq and scRNA-seq, which is competitive with commercial solutions. Second, we developed \texttt{PUMATAC}, a method-agnostic framework for pre-processing and evaluating scATAC-seq datasets, and used it to perform a systematic benchmark of current methods. We open-sourced all protocols and software, opening possibilities for further improvement. HyDrop has since been reproduced in other labs, and derivative products are available commercially. Together, the work touches on all low-level aspects of omics data generation: we optimized reaction chemistries, engineered hardware, wrote software, managed large data structures, and coordinated high-power compute.

\subsection*{\inserted{Development and evaluation of single-cell assays}}

Despite our advancements in open-sourcing methods, microfluidic droplet-based assays still require expertise to implement. Automating flow focusing droplet microfluidics is generally challenging, especially with standard automation hard- and software. The 10x Chromium Connect, an OEM Hamilton Nimbus liquid handler paired with an on-deck Chromium microfluidic controller, offers fully automated droplet generation and library prep, but is serial in nature: one batch of up to 8 samples is processed at a time. This approach thus does not increase sample throughput without human intervention, and is best-suited for small experimental scales. Conversely, Tahoe-100M was generated mostly by one person using 11 well-plates of Parse Bio's Evercode WT Mega kits, which are also amenable to automation since they are plate-native.

Droplet generation itself is also a confounding process that affects the entire sample. If the process fails by the hand of dust, cell or DNA aggregates, air bubbles, chip imperfections, pump controller or gasket failures, etc., the entire sample is affected or lost completely. Manual droplet generation on an Onyx or Chromium followed by partial automation of the library prep is also impractical for many (96+) individual samples due to staggered delivery of samples and general sensitivity to delays. Plate-based single-cell methods can deal with sample multiplexing natively, for example by pre-sorting cells into plates using FACS, freezing the plates, and thawing the plates for processing on demand - a workflow which can be automated end to end. Moreover, the most recent SMART-seq assays have scaled down reagent cost through acoustic nanoliter dispensing and \mbox{\num{1536}} well-plate formats, and combinatorial indexing methods have come closer to 10x Chromium in terms of sensitivity.

Does this mean that droplet-based assays have no practical application for large scale data generation? Of course not. Tahoe-100M's cells were filtered on a minimum of 700 \deleted{umis }\inserted{UMIs }and 250 genes: 10x scRNA-seq undoubtedly still offers higher per-cell quality metrics. Furthermore, since the release of the 10x scATAC-seq v2 kits, which can now reach up to \mbox{\num{100000}} unique fragments per cell, it has become clear that the sensitivity ceiling for droplet-based methods has not been attained by open-source droplet or plate-based scATAC-seq methods. For scRNA-seq, Smart-seq remains the sensitivity gold standard, but has yet to be applied to the same scales as even modest droplet-generated atlases - likely also owing to the high sequencing burden associated. Given only a small slowdown in 10x Genomics' annual kit sales relative to the general capital slump in biotech and research post-ZIRP, it appears that end-users are still willing to make the sensitivity/convenience/cost trade-off in favour of commercial droplet-microfluidics. Thus, how should we further develop the technology? Let us examine outstanding problems and their potential solutions:

\begin{enumerate}[noitemsep]
    \item \textbf{Sample throughput is limited, even on parallel droplet generation chips.} To address this issue, antibody, lipid or oligonucleotide hashes or pre-indexing have been shown to scale throughput by at least 15x, and can scale both samples and absolute number of cells. These multiplexing methods are generally directly compatible with HyDrop's oligonucleotide sequence designs, but none of them are currently compatible with multiomics assays.
    \item \textbf{Emulsion-making is difficult to automate.} PIP-seq, a recent physical sample partitioning development, could offer a more generalizable sample throughput solution. Here, droplet partitioning is achieved by \deleted{vortexing aqueous }\inserted{particle-templated emulsification: an aqueous phase }(cells, enzyme mix \deleted{, bead mix) and lipid phases until a stable emulsion is formed}\inserted{and a hydrogel bead template) is vortexed with a fluorinated oil phase containing surfactant, and each bead templates one droplet of defined volume as the emulsion forms}. This process can be performed in 96-well plates, and can be fully automated! The original PIP-seq article only shows scRNA-seq through this method, and implementing scATAC-seq would pose some challenges: droplet PCR does not pair well with solid or hydrogel beads - the beads must dissolve to prevent emulsion coalescence. HyDrop and 10x achieve this by disulfide bonds which are reduced by contact with the cell and enzyme mix. PIP-seq requires cell and reagent mix to be pre-mixed and vortexed together, so the disulfide approach would not work. However, inDrop has previously used o-nitrobenzyl-mediated UV cleavage to release primers, an idea which could potentially be translated to PIP-ATAC-seq: UV-cleavable beads could be vortexed together with single-cells and enzymes, and the resulting emulsion could then be exposed to UV post-hoc. While this is impractical for manual lab processing, it should not pose a problem for a lights-out automated workflow. Fluent BioSciences, which owns the IP to PIP-seq, and commercialized the assay, was recently acquired by Illumina - promising further development.
    \item \textbf{In-house bead production requires microfluidics and expertise.} After the HyDrop \deleted{Elife }\inserted{eLife }paper was published, we were contacted by several labs who wanted to produce in-house beads, and correspondence revealed that bead production was the largest hurdle to implementation. One common issue was difficulty producing stable emulsions for bead generation, which may possibly be addressed by using PIP-seq-style emulsion generation for bead production as well. Aside from emulsion making, the current barcoding protocol is already fully automatable on existing hardware.
\end{enumerate}

The issues presented above are related mostly to distribution and convenience rather than physical limits of droplet technology. Sensitivity differences between the open-source assay compared to the commercial alternative also play a major role in adoption, and it would thus pay off to tip the balance closer towards the open-source alternative. Since the original publication of HyDrop, an improved version has been developed and open-sourced, offering vastly improved sensitivity as well as more robust bead barcoding protocols. In this updated protocol, bead production was optimized for purity, strongly improving sequencing efficiency and reaching quality parity with 10x. Additional optimisation of extraction and amplification procedures \deleted{are }\inserted{is }likely to further push the sensitivity needle in favor of open-source protocols.

\inserted{Awareness of the sensitivity and specificity biases described here is especially important in scATAC-seq studies compared to scRNA-seq studies, since peak region sets are sample-dependent rather than fixed gene sets. Therefore, high-specificity methods should be prioritised when generating reference peak region sets if public references (e.g. ENCODE ChIP-seq or DNase-seq) are not available a priori. Cost-efficient but lower sensitivity/specificity open-source methods can be used for data generation at scale. Sequencing efficiency of these open-source methods can be improved by sequencing less deeply, since fewer reads will be lost to duplicates, and available resources can be allocated to sampling more cells rather than sequencing more. This strategy will lead to lower per-cell sensitivity, but increasing the total number of cells may be able to partially compensate. For example, both cell type annotation quality (\texttt{seurat} label transfer scores) and DAR strength scaled with higher cell counts (fig. 3a). However, it is likely that scaling data horizontally (by cells only) will generate diminishing returns dependent on minimal baseline per-cell sensitivity required.
}

\subsection*{\inserted{Scaling omics data for predictive modelling}}

Overall, our project feeds into a broader, cross-disciplinary trend of scaling large data explicitly for model training. This trend first began to show results in games like chess and Go, but has recently manifested in more tangible applications such as medical imaging, self-driving cars, and language models. For example, whereas Google's BERT was \deleted{trained on just }\inserted{pre-trained on BookCorpus and English Wikipedia, some }\mbox{\SI{16}{\giga\byte}} of uncompressed text\inserted{, }in 2019, OpenAI's GPT-3 was trained \deleted{on a total of \textasciitilde{\mbox{\SI{700}{\giga\byte}}} of compressed data sourced from }\inserted{just a year later on }\inserted{\mbox{\num{300}}} \inserted{billion tokens drawn from a mixture totalling }\inserted{\mbox{\num{499}}} \inserted{billion tokens: filtered }Common Crawl (\deleted{\mbox{\SI{570}{\giga\byte}}} \deleted{), Wikipedia}\inserted{410B), two internet-based book corpora (12B and 55B)}, WebText2 \deleted{and nearly all written books (}\deleted{\mbox{\SI{120}{\giga\byte}}} \deleted{)just a year later}\inserted{(19B) and Wikipedia (3B)}. Compute was scaled \inserted{by roughly three }orders of magnitude \deleted{more, roughly }\deleted{\mbox{\SI{7.8e20}{\flops}}} \deleted{versus }\inserted{over the same interval, from }\inserted{\mbox{\SI{5.3e20}{\flops}}} \inserted{for BERT-Large to }\mbox{\SI{3.14e23}{\flops}} \inserted{for GPT-3 175B}. The resulting performance improvement, which was primarily the result of a large commitment of data and compute, has allowed GPT-3 to penetrate into consumer markets, eventually becoming the most rapidly adopted technology of all time. Exploring the parameter sensitivity ("scaling laws") of transformer models has shown that their performance continues, and is likely to continue, to scale sublinearly with model dimensions, training dataset size, and compute.

\inserted{Biology is almost certainly at an inflection point in world-total data generation capacity. The cost of sequencing and quantitative sequencing assays has been falling rapidly at predictable rates. What was at the start of this work considered a large atlas of }\inserted{\mbox{\num{100000}}} \inserted{cells can now be routinely generated in a few days, and 10 to 100 million cell datasets are being released. }A question that is much debated today is \textbf{whether the data + compute scaling paradigm will generalize to biological models}. Increasing model and dataset scale has improved performance on several evaluated biological tasks. Gains depend on the task, data composition, and evaluation design. Especially transformer-based protein-language models, notably Meta's ESM family, seem to obey the same scaling laws as their natural language counterparts: between 2020 and 2025, the model family iterated from 650M to 15B and recently 98B parameters, adding generative features.

However, deeper scrutiny reveals more nuance. First, architectural novelty has in some cases outshined pure scaling, especially in regimes where data is scarce. For example, AlphaFold 1 was already trained on a large fraction of available PDB and MSA data, and AlphaFold 3's weights are only \mbox{\SI{1}{\giga\byte}} in size. Introduction of attention with Evoformer blocks and end to end diffusion led to the largest performance jumps. Recent benchmarks have also shown that protein/ligand binding models suffer performance drops when re-trained with more stringent train-test leakage filtering. Performance on ligands that were sparsely represented in the training data was especially weak, suggesting that \textit{"data scarcity remains a major challenge for these methods"}.

Second, smarter data engineering can rival or exceed sheer volume. Nucleotide Transformer improved more on human tasks by adding 850 non-human genomes than by increasing model size from 500M to 2.5B parameters. Evo 2 \deleted{, another transformer model, showed that }\inserted{was released as }a 7B checkpoint trained on \deleted{carefully re-weighted, gene-proximal sequences matches }\inserted{2.4T tokens and }a 40B \deleted{model on most clinical benchmarks while requiring only a tenth }\inserted{checkpoint trained on 9.3T tokens, so the smaller model consumed roughly one twentieth }of the training compute. \deleted{Larger models mainly unlocked generative capabilities: only }\inserted{Both were pretrained with data weighting focused on genic windows. On clinical variant benchmarks the two perform competitively: the 7B and 40B models rank together ahead of unsupervised baselines on coding SNVs, on BRCA1 saturation mutagenesis and on SpliceVarDB, and the 40B pulls clearly ahead mainly on noncoding variants. Scale showed up more visibly in generation quality than in raw capability: both models produced coherent mitochondrial-scale sequence, but tetranucleotide usage of generated sequence agreed more closely with natural genomes for }the 40B \deleted{Evo 2 modelcould synthesize coherent genome-scale DNA sequences}\inserted{model}.

Third, smaller, application-specific models routinely outperform large foundation models in constrained contexts. Bespoke, domain-focused networks trained from scratch on single-cell data such as DeepMEL for melanoma, DeepFlyBrain for \textit{Drosophila}, and DeepZebrafish for embryogenesis, reach higher precision than Enformer and ChromBPNet on enhancer logic while only a few million parameters in size.

\deleted{One modality is palpably absent in the analysis above: perturbation effect prediction.
The }\inserted{Given the high labor and financial cost associated with (single-cell) sequencing assays, and the open question whether the scaling strategy will pay off proportionally, naively launching large omics scaling operations has not been entirely derisked. Considering the nuances above, it might not be the optimal path forward with our current knowledge today, and perhaps focusing on deeper sampling of individual samples might be a better strategy.
}

\inserted{When comparing quantitative omics with DNA language modeling, we must note that the }amount of pure genomic data available is large: \deleted{genBank contains }\deleted{\mbox{\SI{25}{\tera\bp}}} \deleted{from 3.7B annotated sequences in }\inserted{GenBank contains }\inserted{\mbox{\SI{34}{\tera\bp}}} \inserted{from 4.7 billion nucleotide sequences representing }\mbox{\num{581000}} \deleted{species, compared to the \textasciitilde{1.8 million} named eukaryotic species. The training data is thus representative of evolutionary space. }\inserted{formally described species. }Conversely, the quantitative omics space has been sampled much more sparsely\deleted{: GEO currently contains about }\deleted{\mbox{\num{250000}}} \deleted{datasets across 7.5 }\inserted{. At the time of its most recent description, GEO held over }\inserted{\mbox{\num{200000}}} \inserted{studies and 6.5 million samples; an Entrez query in September 2026 returns }\inserted{\mbox{\num{295798}}} \inserted{series across 8.7 }million samples, of which \deleted{around }\deleted{\mbox{\num{380000}}} \deleted{are (sc)RNA-seq, and }\deleted{\mbox{\num{45000}}} \deleted{are (sc)}\inserted{\mbox{\num{144954}}} \inserted{series are expression profiling by high-throughput sequencing, }\inserted{\mbox{\num{50953}}} \inserted{series are genome binding or occupancy profiling, and fewer than }\inserted{\mbox{\num{10000}}} \inserted{series mention }ATAC-seq \inserted{at all}. High-quality perturbation datasets are even fewer in number. This is especially problematic considering that the effect space is also highly combinatorial for both genetic and chemical perturbations. \deleted{Note that in our benchmark we have shown high levels of inefficiency in most omics datasets: gigabytes of \texttt{.fastq}-formated sequencing data }\inserted{Gigabytes of \texttt{.fastq}-formatted sequencing data also }often reduce to sparse \inserted{count }matrices of fewer than \mbox{\SI{50}{\mega\byte}} in size. \inserted{In this regard, quantitative omics might be too sparse yet to derive scaling curves.
}\deleted{Recent work begins to close this data sparsity gap:
Tahoe has delivered }\inserted{\textbf{Assuming a scaling strategy will hold for quantitative omics}, several challenges lie ahead still. Because of the large amount of capital expenditure associated with even modestly-sized omics projects, the main problem is identifying which data types to commit to. Assuming a fixed budget, and a goal of "predict gene regulation in a given disease" (essentially, current virtual cell efforts), several factors affect how we should design a data collection effort:
}

\begin{enumerate}[noitemsep]

\item \inserted{\textbf{Which data type(s) should we collect?} Within sequencing omics assays, we would need to decide on modality (DNA versus accessible chromatin, RNA, Hi-C, methylation, ...), bulk versus single-cell, long versus short reads and so on. Beyond sequencing, there are also dozens of imaging techniques, intracellular and surface protein measurements, functional assays, ... etc. available. There are certainly interaction terms between these different options - for example, ATAC + RNA might be more informative than either alone. We currently do not have a good quantification of "information content per money spent" for various omics techniques. In LLM pre-training, smaller training runs are used to calculate expected information gain and marginal loss reduction, but text acquisition is approximately costless compared to omics data.
}

\item \inserted{\textbf{On which sample?} Omics measurements are generally destructive (although live-cell imaging methods can now monitor selected RNA species). As a result, studies in humans have concentrated around blood, tissue biopsies, and banked samples. While the translational benefit of testing on human samples is directly evident, these samples are not only scarce in absolute terms, but also difficult to obtain due to regulatory and privacy concerns. From a pure information gathering perspective, it is unclear if focusing exclusively on human tissue is the correct call today. For a given budget, substantially deeper sampling of model organisms may prove to be valuable: regulatory programs are conserved across surprisingly distant species. Even simple yeast retains nucleosomes and histone modifications, including methylation, although it lacks cytosine DNA methylation system found in mammals. Cultivating yeast and rearing flies, across many strains and mutants, is more practical than clinical sample collection at scale, and well-controlled genetic backgrounds are a significant benefit. It is thus possible that model organism data at scale could yield models that transfer well to human genomes. Patient-derived organoids could offer a compromise between clinical relevance and experimental scalability. Culture variability remains a limitation, with organoid differentiation efficiency being highly dependent on input iPSC.
}

\item \inserted{\textbf{Investigating which perturbation(s), if any?} Some of the largest current data collection efforts use massively parallel perturbation sampling with single-cell transcriptomic readouts, including both drug and genetic perturbations. However, parallel perturbations have come at the cost of sampling depth and precise control over dosage. Intelligent, iterative perturbation library design might be a source of optimizations.
}

\item \inserted{\textbf{Using which type of automation?} Even if up-front funding were limitless, scaling up data collection will require more than just parallelizing existing personnel and machines. Today, manual research campaigns require different machines, used over various days, and these combinations are generally unique to a project. Moreover, there is certainly a lot of skill involved in collecting and processing samples. Committing to a single scaled data pipeline, where a tissue or other sample is entered one end, and data comes out the other, will require at least the re-designing of existing machines and their close integration with liquid handlers. Humanoid robots might assist during a transitory period, but for large scaling projects, they will most likely not be deployed beyond maintenance work. Similarly to what droplet microfluidics did for single-cell, automated capillary gel-electrophoresis did for DNA sequencing, flow cytometry did for cell sorting, ... we require additional development of specialized monotask machines. Especially in parallel treatment of many samples, more work is needed. For example, a machine that could ingest an arbitrary tissue, and return a highly viable single-cell suspension \textit{reliably} still seems beyond reach for now.
}

\item \inserted{\textbf{Sampling how many cells, at what per-cell depth?} Even the most sensitive plate-based protocols recover under half of a cell's mRNA molecules, and the droplet and split-pool chemistries used for atlas-scale data recover roughly 5-15\% (3.4\% for CEL-seq, 7.1\% for inDrop, 12.8\% for Drop-seq, 6.7\textendash{}8.1\% for 10x). We can thus tune sequencing and assay sampling depth, with number of cells sequenced. For genomic assays, a similar trade-off exists, and single-cell sparsity is even more severe.
}

\end{enumerate}

\inserted{The table places the largest current data-generation efforts on these five dimensions.
}

\inserted{The Human Cell Atlas (HCA), founded in 2016 and still ongoing, prioritizes broader coverage of human cellular diversity. By 2024, the project spanned }100 \inserted{million cells from over }\inserted{\mbox{\num{10000}}} \inserted{individuals, using single-cell and spatial approaches to construct reference maps across tissues. The BRAIN Initiative Cell Atlas Network (BICAN), launched by NIH in 2022 with \$500 million over five years, profiles single-nucleus transcriptome, chromatin accessibility, DNA methylation and spatial transcriptomics from several hundred human brains, and releases data publicly. HuBMAP profiles healthy human organs with multiplexed imaging, spatial transcriptomics and single-nucleus assays.
}

\inserted{Large efforts within Aligning Science Across Parkinson's (ASAP) investigate genetic variation and cell-type-specific regulatory mechanisms in Parkinson's disease. The work reports long-read whole-genome sequencing alongside snATAC-seq profiles from 3.1 million nuclei and snRNA-seq profiles from 1.1 million nuclei, from two brain regions across 190 donors, including 75 individuals with Parkinson's.
}

\inserted{Similar projects are launching in China. BGI-Research and collaborators founded the 10 Billion Cells Alliance in August 2025, proposing reference, disease and perturbation-response atlases using Stereo-cell (cellular imaging + transcriptomics). The Chinese Immune Multi-Omics Atlas, published in 2026, uses single-cell RNA and chromatin accessibility measurements from more than 10 million circulating immune cells across 428 adults.
}

\inserted{CIMA's OneK1K profiled 1.27 million blood cells from 982 people and mapped cell-type-specific eQTLs to autoimmune risk loci. Its follow-up project, TenK10K, launched in 2024 by Garvan with Illumina and 10x, targets 50 million cells and genomes from }\inserted{\mbox{\num{10000}}} \inserted{people.
}

\inserted{A number of startups have also oriented explicitly on a "data-first, model next" path to revenue. Noetik, founded in 2022, generates microarrays from clinical tissues, and profiles them using H\&E histology, spatial transcriptomics, spatial protein measurements and DNA sequencing. Their private model is trained on paired histology and spatial transcriptomic data, and predicts spatial gene expression from H\&E alone. Recursion, another startup, also focuses heavily on imaging, using immunofluorescence plus chemical perturbations. Their public RxRx3 resource includes more than }\inserted{\mbox{\num{17000}}} \inserted{gene knockouts and }\inserted{\mbox{\num{1600}}} \inserted{compounds. Xaira's X-Atlas/Orion (2025) applied genome-wide CRISPRi to two cell lines and sequenced eight million cells to more than }\inserted{\mbox{\num{16000}}} \inserted{UMIs each. X-Atlas/Pisces extends this to 25.6 million cells across seven cellular contexts, currently the largest Perturb-seq dataset. Owkin's MOSAIC, a \$50 million consortium with five European and US cancer centres, profiles }\inserted{\mbox{\num{7000}}} \inserted{tumour samples across seven cancer types with GeoMx, CosMx, Visium and Chromium. Sixty patients are open access, the rest is held private.
}

\inserted{IGVF (NHGRI, 2021-2026) maps chromatin accessibility and transcriptome with single-nucleus multiome in cell lines, iPSC-derived and primary cells, perturbs variants, elements and genes with CRISPR screens, MPRAs and saturation mutagenesis, and trains models of variant effect on the result. It is the only consortium that pairs perturbation with a regulatory readout, at $10^4$-$10^5$ cells per experiment.
}

\inserted{Tahoe-100M contains 100 }million \deleted{sci-derived }single-cell transcriptomes of 50 \deleted{diverse cancer lines perturbed with 1 to 100 small-molecule conditions (\textasciitilde}\deleted{\mbox{\num{17800}}} \deleted{unique }\inserted{cancer cell lines treated with 379 small molecules: }\inserted{\mbox{\num{1135}}} \inserted{drug-dose conditions, }\inserted{\mbox{\num{52886}}} line-drug \deleted{pairs). \texttt{State}, an accompanying transformer model , contrasts this perturbation atlas with 167M observational cells sourced from public data to predict the effect of perturbations, exceeding mean-baselines for zero-shot ranking of drug effect sizes in unseen cell types by > 15\%. Retraining \texttt{State} on downsampled scBaseCamp could also provide the first quantitative scaling laws for omics modeling}\inserted{conditions, median }\inserted{\mbox{\num{1287}}} \inserted{cells per condition. After Tahoe's initial proof of concept, Arc, Biohub and Tahoe announced a joint effort of more than 120 million cells spanning }\inserted{\mbox{\num{225000}}} \inserted{drug-patient interactions, with more than four times the perturbation richness of Tahoe-100M}.

\deleted{Given the high labor and financial cost associated with (single-cell)sequencing assays}\inserted{The Billion Cells Project, launched between CZI, 10x and Ultima Genomics in 2025, targets one billion single-cell profiles including genetic perturbations, in mouse, zebrafish and primary human cell models. Illumina's competing Billion Cell Atlas, in partnership with Eli Lilly, AstraZeneca and Merck, now aims to measure CRISPR-perturbation responses using scRNA-seq in 200 disease cell lines. The work relies heavily on the PIP-seq IP that Illumina re-packaged into its Illumina Single Cell 3' RNA Prep kit. As of July 2026}, \inserted{350 million cells were completed.

\textit{[Table placeholder: large-scale omics data-generation efforts mapped onto the five design dimensions.]}
}

\deleted{and the open question whether the scaling strategy will pay off proportionally, naively launching extreme scaling operations has not been entirely derisked. Considering the nuances above, it is unlikely to be the optimal path forward with our current knowledge. However, we are almost certainly at an inflection point in world-total data generation capacity. The cost of sequencing and quantitative sequencing assays has been falling rapidly at predictable rates. What was at the start of this work considered a large atlas of \mbox{\num{100000}} cells can now be routinely generated in a few days, and 10 to 100million cell datasets are being released. }

\inserted{We can see trends emerge in current mega-projects. Generally speaking, academic consortia more often work with native tissues, measure several modalities per donor, and publish their findings and data as public reference resources. Industry opts for a broad perturbation approach, often in cell lines, a decision which stems from the industry's aim to ultimately predict the effects of future therapeutic interventions. Perturbation efforts read out either the transcriptome or images, none reads out a second sequencing modality. Chromatin accessibility appears only in the observational projects (ASAP, CIMA). Perturb-ATAC, CRISPR-sciATAC and Spear-ATAC established this is technically possible, but have so far only been demonstrated at scales of up to 100,000 cells. Mouse and zebrafish appear just once, yeast and \textit{Drosophila} not at all. Arc, Tahoe and Roche discuss future plans in organoids, but as of 2026, no data has been published yet. The largest cell counts come from plate-native, internal strategies: SPLiT-seq (Tahoe), PIP-seq (Illumina) and fixed-cell batching (X-Atlas). The observational consortia scale instead by distributing conventional droplet workflows across many laboratories, and as a result achieve better quality depth per sample.
}

\inserted{Together, this cohort of collection projects suggests both industry and academic confidence in omics data scaling strategies has grown significantly. }It is also likely that the process of data scaling will self-reinforce. Better understanding of proteins and gene regulation will lead to development of better molecular tools such as synthetic enzymes, optimal expression systems, and better perturbation delivery methods. Improved models will also improve experimental design through more accurate hypothesis (and thus informative data) generation, and may even generate synthetic data for finetuning. Advances in lab automation coupled with automated protocol generation through LLM support (PyLabRobot) will further accelerate this process.
\deleted{Finally, several of the models and datasets discussed here have been produced by privately-funded research institutions (Google DeepMind, Meta, CZI, Calico, Arc, Tahoe), and several automated omics assay providers (Plasmidsaurus, Adaptyv Bio) are emerging. There is thus mounting evidence that industry confidence in omics data scaling strategies is growing. Whereas today's biological models gather mostly all available data to produce a state-of-the-art result, a "GPT moment" for biology could spur a data economy in which model trainers generate a demand for specific data. This switch from a supply-driven to a demand-driven economy has effectively already happened in consumer AI, as thousands of people are now employed as data labelers for various data modalities, and new methods for embodied and physical data generation are being commercialized}\inserted{It remains yet to be seen which strategy will contribute positively to drug discovery and therapy development. If all goes well, the emerging data ecosystem might contain several complementary types of resource: a mixture of broad reference atlases, characterized patient cohorts, and large private perturbation models}. As model predictive performance increases, downstream medicine and consumer products will further stimulate capital attention and \deleted{thus }model scale. In that sense, the biotech and pharmaceutical industry may soon become a data economy where information becomes the primary currency of discovery.

\endgroup
\end{document}
